\section{A common finite count comparison}\label{ext:finite-counts}
The finite argument follows the number of positive summands. Conditional on
that number, the sum is a binomial lattice point plus a centered residual.
The residual can move mass across a threshold in both directions. Keeping
its outgoing and incoming contributions together is essential near a
binomial jump. Everything in this section concerns one common marginal;
no finite-atom assumption is made.

\subsection{Exact conditioning and the common moments}
For a standardized law put
\[
 p=\Pp(X>0),\quad q=1-p,\quad v=pq,\quad m=\E X_+,
 \quad h=v/m.
\]
Centering and positive variance imply $0<p<1$ and $m>0$. With
$B=\ind_{\{X>0\}}$, define $U=hX-B+p$. Then
$\E(U\mid B)=0$, and
\[
 hX=B-p+U,\qquad h^2=v+s,\qquad
 s=q w_0+p w_1,\quad w_i=\E(U^2\mid B=i).
\]
The atom at zero belongs to $B=0$. If $s>0$, let
$\alpha=pw_1/s$ and $\lambda=ns$; thus
$w_0=\lambda(1-\alpha)/(nq)$ and $w_1=\lambda\alpha/(np)$.
When $s=0$, both conditional residuals vanish and these formulas hold
with any $\alpha\in[0,1]$.

Write $W_j$ for a sum of $n-j$ independent copies of $U\mid B=0$
and $j$ independent copies of $U\mid B=1$, all independent. For
$b_j=\binom njp^jq^{n-j}$ and $z=np+hH=k+u$,
\begin{align}
 \Pp(S_n\le H)&=\sum_{j=0}^n b_j\Pp(W_j\le k+u-j),
 \label{ext:count-identity}\\
 \Pp(S_n\le H)&\le F_{\mathrm{Bin}(n,p)}(k)
 -b_k\Pp(W_k>u)+\sum_{j>k}b_j\Pp(W_j\le u-(j-k)).
 \label{ext:count-bound}
\end{align}
The second line discards nonnegative outgoing probabilities at $j<k$.
The outgoing event is strict and the incoming event is closed. In particular,
an atom exactly on the threshold is retained on the CDF side.
Also
\[
 V_j:=\E W_j^2=(n-j)w_0+jw_1
 =\lambda\left(1+\frac{(j-np)(\alpha-p)}{npq}\right).
\]

Suppose support acquisition gives
$U_0\in[-O_0,a_0]$ and $U_1\in[-a_1,O_1]$, with nonnegative
endpoints. For a centered $T\in[-L,R]$ of variance $w$,
\begin{equation}\label{ext:conditional-moments}
 w\le LR,\quad w^2/L-Lw\le\E T^3\le Rw-w^2/R,
 \quad \E T^4\le(R-L)\E T^3+LRw.
\end{equation}
The first inequality integrates $(T+L)(R-T)\ge0$. The two cubic
bounds integrate $(T+L)(T-w/L)^2\ge0$ and
$(R-T)(T+w/R)^2\ge0$. Multiplying the first polynomial by $T^2$
gives the fourth-moment bound. If $L=0$ or $R=0$, centering forces
$T=0$; the divisions are replaced by this deterministic case.

Put $\mu_i=\E U_i^3$. Expanding on the two actual sign classes gives
\begin{align}
 \rho h^3&=v(1-2v)+3v(w_0+w_1)-q\mu_0+p\mu_1,
 \label{ext:conditional-rho}\\
 \gamma h^3&=v(1-2p)+3v(w_1-w_0)+q\mu_0+p\mu_1,\\
 \beta h^2&=v(1-2p)+s(2\alpha-1).
\end{align}
Here $\gamma=\E X^3$ and $\beta=\E[X|X|]$. The same
$(\mu_0,\mu_1)$ must be used in these identities and in every probability
bound below. Optimizing each term at a different pair would lose the
constraint that all terms arise from one law.

\subsection{Outgoing and incoming probability bounds}
A useful outgoing minorant is
\[
 G_a(x)=x(x+a)^2(2a-x),\qquad a>0.
\]
It is nonpositive outside $[0,2a]$, and its maximum is
$K a^4$, where $K=9/4+3\sqrt3/2$. Consequently
\[
 \Pp(W_k>u)\ge \frac{\E G_a(W_k-u)}{K a^4}.
\]
This is a global polynomial inequality, including $W_k=u$.
Its expectation uses
\[
 \E W_k^3=(n-k)\mu_0+k\mu_1,
\]
\[
 \E W_k^4=(n-k)\E U_0^4+k\E U_1^4
 +3\{V_k^2-(n-k)w_0^2-kw_1^2\}.
\]
Thus \eqref{ext:conditional-moments} makes the outgoing bound affine in
the two common cubic moments. Negative lower bounds remain valid; they
are not mistaken for negative probabilities.

For an incoming count $j>k$ and phase $u\le1/2$, put $d=j-k-u>0$.
Cantelli gives
\[
 \Pp(W_j\le-d)\le\frac{V_j}{V_j+d^2}.
\]
For completeness, for a centered $Y$ of variance $V$ and $a>0$,
$\Pp(Y\ge d)\le\E(Y+a)^2/(d+a)^2=(V+a^2)/(d+a)^2$.
Choosing $a=V/d$ proves the formula; $V=0$ is deterministic.
Exponential bounds can be sharper. If $\E Y=0$, $Y\le B$, $B>0$,
and $\E Y^2\le w$, then for $t\ge0$,
\begin{equation}\label{ext:one-sided-mgf}
 \E e^{tY}\le
 M(w;B,t):=\frac{B^2e^{-tw/B}+we^{tB}}{B^2+w}.
\end{equation}
Indeed, the quadratic interpolating $e^{tx}$ at $-w/B$ with multiplicity
two and at $B$ majorizes it on $x\le B$: the Hermite remainder has
sign $(x+w/B)^2(x-B)\le0$. Its quadratic coefficient is nonnegative.
The centered two-point law on $\{-w/B,B\}$ with variance $w$ evaluates
its expectation exactly. This proves the formula also with an upper
variance bound. At zero variance use $Y=0$.
Differentiation shows that $M$ is increasing and concave in $w$:
with $x=t(B^2+w)/B$,
\[
 M''(w)=\frac{B^2e^{-tw/B}}{(B^2+w)^3}
          (x^2+2x+2-2e^x)\le0.
\]
Also
\[
 M'(w)=\frac{B^2}{(B^2+w)^2}
 \{e^{tB}-e^{-tw/B}(1+x)\}\ge0,
\]
since $e^x\ge1+x$. Its tangents therefore give affine upper bounds throughout the variance
half-line. Applying \eqref{ext:one-sided-mgf} to $-U_0,-U_1$ and
multiplying the corresponding factors bounds $\E e^{-tW_j}$.
Chernoff then gives $e^{-td}\E e^{-tW_j}$ as an incoming upper bound.

One can also retain the common cubic moments. On $[A,B]$, $A<0<B$,
the function
\[
 r_t(x)=\frac{e^{tx}-1-tx}{x^2}
       =t^2\int_0^1(1-r)e^{trx}\,dr
\]
is increasing and convex. Its chord gives
$e^{tx}\le1+tx+c_2x^2+c_3x^3$, with $c_2,c_3\ge0$.
Applied to $-U_i$, choose $C_2\ge c_2$ and
$0\le C_{3,-}\le c_3\le C_{3,+}$, and put
\[
 Q_i(\mu_i)=\max\{0,1+C_2w_{i,+}-C_{3,-}\mu_i,
                     1+C_2w_{i,+}-C_{3,+}\mu_i\}.
\]
This is a nonnegative convex upper bound for its MGF. The maximum with zero
also makes it nonnegative at infeasible vertices of a relaxed moment polygon.
For fixed positive pivots $P_0,P_1$, weighted AM--GM yields
\[
 Q_0^{n-j}Q_1^j\le P_0^{n-j}P_1^j
 \left(\frac{n-j}{n}\frac{Q_0}{P_0}
       +\frac{j}{n}\frac{Q_1}{P_1}\right)^n.
\]
The right side is convex in $(\mu_0,\mu_1)$. A fixed exponential rate
and its phase chord produce a convex common-moment bound on an entire
closed phase interval. Choose a strategy on that whole interval and
moment polygon before comparing different strategies; a pointwise
minimum of convex functions need not be convex.

\subsection{Normalization and a finite closed cover}
Let $\mathcal M$ be a compact polygon containing the actual common
cubic moments, obtained from proved moment and defect restrictions.
Suppose $F$ is convex on $\mathcal M$ and $L$ is affine and strictly
positive there. Then
\begin{equation}\label{ext:normalized-vertex}
 \frac{F(x)}{L(x)}\le
 \max_{v\text{ vertex of }\mathcal M}\frac{F(v)}{L(v)}.
\end{equation}
To see this, write $x=\sum_v\theta_vv$. Convexity bounds the numerator
by $\sum_v\theta_vF(v)$; affinity identifies the denominator with
$\sum_v\theta_vL(v)$. The ratio is therefore bounded by the largest
vertex ratio. Segments and single points are included. If the lower
bound for $\rho$ is $\max\{r_0,L\}$, split $\mathcal M$ along
$L=r_0$ and retain both closed pieces.

\begin{proposition}[Whole-cell count certificate]\label{ext:count-cell}
Fix a closed box of acquired parameters, a count $k$, and a closed
phase interval. Let $B_k$ bound the binomial CDF above, $b_k^-$ the own
mass below, and $b_j^+$ the incoming masses above. Let $L_k$ be a
proved outgoing lower bound, $A_j\ge0$ proved incoming upper bounds,
and $N$ a lower bound for $\Phi((k+u-np)/(\sqrt n h))$.
Suppose the combination
\[
 B_k-b_k^-L_k+\sum_{j>k}b_j^+A_j-N
\]
is bounded on the whole phase interval by convex Bernstein control
functions $F_i(\mu_0,\mu_1)$. Let $r>0$ be the affine lower bound for
$\rho$ on each closed normalization piece and let $n_+\ge\sqrt n$.
Then the positive normalized CDF error is at most
\[
 n_+\max_{i,\,v\text{ vertex}}
       \frac{\max\{0,F_i(v)\}}{r(v)}.
\]
\end{proposition}
\begin{proof}
Insert the probability bounds into \eqref{ext:count-bound}. For the outgoing
term, $b_k\Pp(W_k>u)\ge b_k^-\Pp(W_k>u)\ge b_k^-L_k$,
even if $L_k<0$. The nonnegativity of the incoming bounds justifies the
upper mass substitutions. The Bernstein basis is nonnegative and sums
to one. Take the positive part and apply \eqref{ext:normalized-vertex};
the positive part of a convex function is convex. Finally enlarge
$\sqrt n$ to $n_+$.
\end{proof}

The acquired $p,h,H$ bounds must first put every actual threshold into
one of the original roots $k=0,\ldots,n$, $u\in[-1/2,1/2]$.
If they do not, prove the remaining phase tails separately. A tied
nearest count belongs to both closed roots. Apply
Proposition~\ref{ext:closed-cover} only after every leaf is either a
proved necessary-condition contradiction or satisfies the whole-cell
score bound. The finite tree then covers infinitely many parameters.
Its leaves never assert that every feasible moment tuple is an iid law;
they need only contain all actual laws.

\subsection{The predecessor input and index transfers}
At a ratio maximizer from Corollary~\ref{ext:ratio-attainment},
$C=C_n>\CE$ and $C_{n-1}\le\CE$ imply
\[
 n\left(\frac{\sqrt{n-1}\Delta_{n-1}(P)}{\rho(P)}-C\right)_+=0.
\]
This supplies the zero predecessor allowance for the same law. It is
not a monotonicity assertion about the sequence $C_n$.

When a resource calculation is transferred from a floor $N$ to an actual
index $m\ge N$, the invariant is
\[
 \mathcal R_{\rm actual}\le r,\qquad mD_{\rm actual}\le Nd.
\]
For example, if the same-law defect inequality is
$D\ge b+g d_{\rm in}+(N/m)h_N\lambda_m$ with nonnegative
$b,g,h_N$, multiplying by $m$ and dividing by $N$ gives
$(m/N)(b+g d_{\rm in})+h_N\lambda_m\le d$.
Since $m/N\ge1$ and $b+g d_{\rm in}\ge0$, this implies
$b+g d_{\rm in}+h_N\lambda_m\le d$.
The noise is $\lambda_m=ms$, and cannot be replaced by $Ns$.
Every downstream bound must retain its original index interval.
In particular, $I/m\le\max\{I/40,I/42\}$ for a signed $I$ is a
valid floor transfer only on $40\le m\le42$.
